\documentclass{article}

\usepackage{xcolor}
\usepackage[utf8]{inputenc}
\usepackage{amsmath}
\usepackage{enumerate}
\usepackage{microtype}
\usepackage[hyphens]{url}
\usepackage{hyperref}
\usepackage[top=1.5cm]{geometry} % Page margins

\title{Rúbrica Proyecto Final - Plan de Estudios}
\author{Programación III \\ Facultad de Ingenierías - Universidad Tecnológica de Pereira \\ Ingeniería de Sistemas y Computación}
\date{}

\begin{document}
\maketitle
\thispagestyle{empty}

\section{Objetivo}

Para este proyecto se requiere elaborar un modelo de la malla curricular para el programa de Ingeniería de Sistemas y Computación de la Universidad Tecnológica de Pereira que permita obtener un plan de estudios adaptado a las condiciones requeridas por el usuario.

\section{Requerimientos}

\subsection{Parte I. Modelamiento mediante restricciones}
Modelar el plan de estudios del programa de Ingeniería de Sistemas y Computación teniendo en cuenta el dominio de las variables de decisión y las restricciones asociadas al mismo:

\begin{itemize}
    \item Conjunto total de asignaturas.
    \item Número de créditos asociado a cada asignatura.
    \item Relaciones de prerrequisito entre asignaturas.
    \item Relaciones de correquisito entre asignaturas.
    \item Restricciones relacionadas con la carga académica por semestre.
\end{itemize}


A partir de este modelo, obtener una solución global que construya un plan de matrícula por semestre, cumpliendo todas las restricciones establecidas y permitiendo que un estudiante complete el programa académico en un máximo de 12 semestres.

Debe analizar los resultados obtenidos y reportar hallazgos y conclusiones.

\subsection{Parte II. Agentes}

Teniendo en cuenta el modelo anteriormente definido, debe implementar un sistema de toma de decisiones basado en agentes. Dos ejemplos de implementación de agentes, según lo discutido en clase son:

\begin{enumerate}
    \item \textbf{Agente "Alumno esforzado":} Selecciona las asignaturas disponibles con mayor número de créditos asociados en orden descendente, respetando las restricciones académicas.
    \item \textbf{Agente "Alumno ocupado":} Selecciona asignaturas hasta completar un máximo de 15 créditos por semestre, sin priorizar el número de créditos de cada asignatura.
\end{enumerate}

Adicionalmente, se debe implementar un agente o una sociedad de agentes propia que permita hallar una solución válida para el problema modelado.

\textbf{Requerimiento de Sociedad de agentes:}  Debe incluir mecanismos de interacción, cooperación o competencia entre 2 o más agentes para construir conjuntamente una propuesta válida de plan académico. \newline


El proyecto puede desarrollarse en cualquier lenguaje de programación. Se recomienda el uso de MiniZinc integrado con Python o SWI-Prolog.

\section{Entregables}
Se debe presentar una sustentación del proyecto, así como el código fuente creado para la resolución de los requerimientos definidos.
\end{document}

